\documentclass[10pt,a4paper]{article}
% Compile with pdflatex (e.g. on Overleaf). results.tex is generated by `python -m src.build_docs`.
\usepackage[utf8]{inputenc}
\usepackage[T1]{fontenc}
\usepackage{lmodern}
\usepackage[a4paper,margin=1.8cm]{geometry}
\usepackage{amsmath,amssymb,bm}
\usepackage{graphicx}
\usepackage{booktabs}
\usepackage{xcolor}
\usepackage{caption}
\usepackage{subcaption}
\usepackage{enumitem}
\usepackage[hidelinks]{hyperref}
\usepackage{titlesec}

\captionsetup{font=small,labelfont=bf}
\setlist{nosep,leftmargin=1.2em}
\titlespacing*{\section}{0pt}{1.0ex}{0.6ex}
\titlespacing*{\paragraph}{0pt}{0.6ex}{0.8em}
\setlength{\parskip}{0.35em}
\setlength{\parindent}{0pt}

\newcommand{\todo}[1]{\textcolor{red}{\textbf{[TO FILL AFTER RUN#1]}}}
% AUTO-GENERATED by `python -m src.build_docs` from results/ -- do not edit by hand.
\newcommand{\figdir}{../figures/}
\newcommand{\NTrainRaw}{161,297}
\newcommand{\NTestRaw}{53,766}
\newcommand{\NTrainKept}{146,941}
\newcommand{\NTestKept}{48,937}
\newcommand{\Dropped}{19,185}
\newcommand{\PosTrain}{72.7\%}
\newcommand{\PosTest}{72.4\%}
\newcommand{\LenMean}{94}
\newcommand{\LenMedian}{93}
\newcommand{\LenPNF}{162}
\newcommand{\Coverage}{99.8\%}
\newcommand{\NCond}{898}
\newcommand{\NDrugs}{3,617}
\newcommand{\NTrain}{40,000}
\newcommand{\NVal}{5,000}
\newcommand{\NTest}{10,000}
\newcommand{\Vocab}{18,664}
\newcommand{\HparamFindings}{Relative to the baseline (macro-F1 0.8561): \emph{learning rate}: Learning rate 5e-3 gives 0.8561 (+0.0000), best epoch 5, 48.4 s/epoch. \emph{optimizer}: SGD (momentum 0.9, lr 0.05) gives 0.8348 (-0.0213), best epoch 8, 48.6 s/epoch. \emph{batch size}: Batch size 64 gives 0.8552 (-0.0009), best epoch 8, 63.5 s/epoch. \emph{batch size}: Batch size 256 gives 0.8511 (-0.0050), best epoch 8, 42.0 s/epoch. \emph{initialization}: Orthogonal recurrent init gives 0.8467 (-0.0093), best epoch 4, 48.5 s/epoch. The best setting in this study is \textbf{Learning rate 5e-3} (0.8561). The baseline (Adam, lr 1e-3) is already a good setting. SGD with momentum is the only clearly worse choice (-0.021 macro-F1), and it is under-fitted rather than overfitted: after 10 epochs its training loss is still about 0.20 (vs 0.02 with Adam) and its validation loss never turns upward, so it converges more slowly within the same epoch budget. A 5x larger Adam learning rate reaches the same macro-F1 three epochs earlier (best epoch 5 vs 8). Batch size mainly trades time for stability: batch 64 matches the baseline but is about 30\% slower per epoch, batch 256 is the fastest but loses 0.005. Orthogonal recurrent initialization brought no gain on the test set (-0.009) and early-stopped after 7 epochs; with reviews of about 100 tokens and Adam, its benefit for gradient flow over long sequences is not needed.}
\newcommand{\ArchFindings}{The vanilla baseline reaches a test macro-F1 of 0.8561 (accuracy 0.8872). CIFG-LSTM reaches a test macro-F1 of 0.8614 (+0.0053 vs vanilla) with 25.0\% fewer recurrent parameters and 41.1 s/epoch vs 48.6 s for vanilla; its best validation epoch is 6. LSTMP (P=128) reaches a test macro-F1 of 0.8526 (-0.0035 vs vanilla) with 24.9\% fewer recurrent parameters and 51.2 s/epoch vs 48.6 s for vanilla; its best validation epoch is 8. All three cells land within 0.01 macro-F1 of each other, which a single seed cannot reliably separate, so the main finding is that neither efficiency variant loses measurable accuracy: CIFG removes the input gate (25\% fewer recurrent parameters, about 15\% faster epochs) and even scores slightly higher, and LSTMP keeps a 256-d cell with a 128-d recurrent state for the same budget at a small cost (-0.0035). The validation curves show the same pattern for every cell: validation loss is lowest around epoch 4 and then rises while training loss falls to about 0.02 and validation macro-F1 stays on a plateau between 0.84 and 0.87. The models overfit and become over-confident after a few epochs; selecting the checkpoint on macro-F1 rather than on loss keeps the better-ranked epochs.}
\newcommand{\ProjFindings}{\textbf{Projection size (bonus).} Halving the projection from P=128 to P=64 reduces the recurrent parameters from 296,064 to 214,080 (27.7\% fewer) and changes test macro-F1 from 0.8526 to 0.8493 (-0.0033).}
\newcommand{\BestSummary}{Selecting on \textbf{validation} macro-F1 across all runs, the final model is \textbf{cifg} (cell \texttt{cifg}, 2,685,186 trainable parameters, val macro-F1 0.8684). On the 10,000-review test set it obtains accuracy \textbf{0.8893}, macro-F1 \textbf{0.8614}, macro precision 0.8616 and macro recall 0.8612.}
\newcommand{\CmSentence}{Recall is 79.8\% on negative reviews and 92.4\% on positive reviews: 556 negative reviews are predicted positive and 551 positive reviews are predicted negative. Errors are concentrated on the minority class: the two error types are almost equally frequent in absolute terms (556 vs 551), but negatives are only about 28\% of the test set, so recall is about 12 points lower on negative reviews than on positive ones. A class-weighted loss or a decision threshold tuned on the validation set (Section 6) directly targets this gap.}
\newcommand{\AbstractResult}{The final model, selected on validation macro-F1, is \textbf{cifg}, which reaches a test accuracy of 0.889 and a test macro-F1 of 0.861.}
\newcommand{\ParamsTable}{%
\begin{tabular}{lrrr}
\toprule
\textbf{Cell} & \textbf{Recurrent params} & \textbf{Total params (incl. embedding + classifier)} & \textbf{vs vanilla (recurrent)} \\
\midrule
Vanilla LSTM & 394,240 & 2,783,746 & +0.0\% \\
CIFG-LSTM & 295,680 & 2,685,186 & -25.0\% \\
LSTMP (P=128) & 296,064 & 2,685,314 & -24.9\% \\
LSTMP (P=64) & 214,080 & 2,603,202 & -45.7\% \\
\bottomrule
\end{tabular}}
\newcommand{\ArchTable}{%
\resizebox{\linewidth}{!}{%
\begin{tabular}{lrrrrrrrr}
\toprule
\textbf{Model} & \textbf{Trainable params} & \textbf{Recurrent params} & \textbf{Test acc.} & \textbf{Test macro-F1} & \textbf{Precision} & \textbf{Recall} & \textbf{Best epoch} & \textbf{Time/epoch (s)} \\
\midrule
Vanilla LSTM & 2,783,746 & 394,240 & 0.8872 & 0.8561 & 0.8635 & 0.8495 & 8 & 48.6 \\
CIFG-LSTM & 2,685,186 & 295,680 & 0.8893 & 0.8614 & 0.8616 & 0.8612 & 6 & 41.1 \\
LSTMP (P=128) & 2,685,314 & 296,064 & 0.8806 & 0.8526 & 0.8481 & 0.8575 & 8 & 51.2 \\
\bottomrule
\end{tabular}}}
\newcommand{\HparamTable}{%
\resizebox{\linewidth}{!}{%
\begin{tabular}{lrrrrrrrr}
\toprule
\textbf{Vanilla LSTM variant} & \textbf{Trainable params} & \textbf{Recurrent params} & \textbf{Test acc.} & \textbf{Test macro-F1} & \textbf{Precision} & \textbf{Recall} & \textbf{Best epoch} & \textbf{Time/epoch (s)} \\
\midrule
Baseline (Adam, lr 1e-3, bs 128, Xavier) & 2,783,746 & 394,240 & 0.8872 & 0.8561 & 0.8635 & 0.8495 & 8 & 48.6 \\
Learning rate 5e-3 & 2,783,746 & 394,240 & 0.8857 & 0.8561 & 0.8582 & 0.8541 & 5 & 48.4 \\
SGD (momentum 0.9, lr 0.05) & 2,783,746 & 394,240 & 0.8632 & 0.8348 & 0.8249 & 0.8477 & 8 & 48.6 \\
Batch size 64 & 2,783,746 & 394,240 & 0.8838 & 0.8552 & 0.8537 & 0.8568 & 8 & 63.5 \\
Batch size 256 & 2,783,746 & 394,240 & 0.8836 & 0.8511 & 0.8596 & 0.8438 & 8 & 42.0 \\
Orthogonal recurrent init & 2,783,746 & 394,240 & 0.8789 & 0.8467 & 0.8508 & 0.8430 & 4 & 48.5 \\
\bottomrule
\end{tabular}}}


% \fig{file}{width}{caption}: shows a placeholder box until the figure exists
\newcommand{\figbox}[2]{%
  \IfFileExists{\figdir #1}{\includegraphics[width=#2]{\figdir #1}}%
  {\fbox{\parbox[c][3cm][c]{#2}{\centering\todo{}\\ \small\texttt{#1}}}}}

\title{\vspace{-1.2cm}\textbf{Sentiment Classification of Drug Reviews\\with Hand-Implemented LSTM Variants}\\[0.3em]
\large Advanced Natural Language Processing --- Lab Project on Recurrent Neural Networks}
\author{Mohamed Amine Lamrid \and Badr Misraoui\\[0.2em]
\small Lab group E \quad$\cdot$\quad University year 2026/2027\\
\small Code: \url{https://github.com/amine2003lamrid-hue/drug-review-lstm}}
\date{}

\begin{document}
\maketitle
\vspace{-1.0cm}

\begin{abstract}
\noindent We classify patient drug reviews from Drugs.com as positive or negative with recurrent networks whose
cells are implemented from scratch in PyTorch using only \texttt{nn.Linear}. A vanilla LSTM baseline is compared with
two efficiency-oriented variants, the Coupled Input--Forget Gate LSTM (CIFG) and the LSTM with recurrent projection
(LSTMP), under an identical protocol, and the baseline's sensitivity to learning rate, optimizer, batch size and weight
initialization is studied one factor at a time. \AbstractResult
\end{abstract}

%-----------------------------------------------------------------------------------------------
\section{Goal of the task}
The UCI Drug Review dataset contains patient reviews of medications rated from 1 to 10. We frame it as
\textbf{binary sentiment classification}: a review is \emph{positive} if its rating is $\geq 7$ and \emph{negative} if
it is $\leq 4$; neutral reviews (5--6) are discarded. The study has three goals: (i) build a vanilla LSTM baseline;
(ii) test whether two non-vanilla cells, \textbf{CIFG} and \textbf{LSTMP}, can match it with fewer parameters;
(iii) measure how optimization hyperparameters affect the baseline. All cells are written with \texttt{nn.Linear} as
the only linear operation; no built-in recurrent module (\texttt{nn.LSTM}, \texttt{nn.LSTMCell}, \texttt{nn.RNN}) is
used, which a unit test verifies automatically.

%-----------------------------------------------------------------------------------------------
\section{Exploratory data analysis}
The official files contain \NTrainRaw{} training and \NTestRaw{} test reviews. Dropping the \Dropped{} neutral reviews
leaves \NTrainKept{} training and \NTestKept{} test reviews. Ratings are strongly polarized (10 and 1 dominate,
Fig.~\ref{fig:eda}a), so the binary task is well defined, but the classes are \textbf{imbalanced}: \PosTrain{} of the
training reviews are positive (\PosTest{} in test, Fig.~\ref{fig:eda}b). A constant ``positive'' classifier would
therefore already reach about \PosTest{} accuracy, so we use \textbf{macro-F1} for model selection and as the headline
metric.

\begin{figure}[h]
\centering
\begin{subfigure}{0.49\linewidth}\centering\figbox{rating_distribution.png}{\linewidth}\caption{Ratings (orange: negative, gray: dropped, blue: positive)}\end{subfigure}\hfill
\begin{subfigure}{0.49\linewidth}\centering\figbox{class_balance.png}{\linewidth}\caption{Class balance after binarization}\end{subfigure}\\[0.3em]
\begin{subfigure}{0.49\linewidth}\centering\figbox{review_length.png}{\linewidth}\caption{Review length in tokens (train)}\end{subfigure}\hfill
\begin{subfigure}{0.49\linewidth}\centering\figbox{top_conditions.png}{0.9\linewidth}\caption{Most reviewed conditions}\end{subfigure}
\caption{Exploratory data analysis of the Drug Review dataset.}
\label{fig:eda}
\end{figure}

\paragraph{Preprocessing.} Reviews are HTML-unescaped (the raw text contains entities such as \texttt{\&\#039;}),
lowercased, stripped of quotes and tokenized with a regular expression keeping words, contractions and the punctuation
marks \texttt{! ? . ,}. The mean review length is \LenMean{} tokens (median \LenMedian{}, 95th percentile \LenPNF{}),
with a sharp drop around 160 tokens that reflects a length limit on the website (Fig.~\ref{fig:eda}c). A maximum
length of \textbf{200 tokens} keeps \Coverage{} of reviews untruncated. The reviews cover \NCond{} conditions and
\NDrugs{} drugs, led by birth control, depression and pain (Fig.~\ref{fig:eda}d).

\paragraph{Splits.} The official test file is the test set; the official training file is split 90/10 (stratified,
seed 42) into training and validation sets. To fit free-GPU budgets we draw \textbf{stratified} subsamples of
\NTrain{} / \NVal{} / \NTest{} reviews, which preserve the original class ratio and so introduce no artificial
imbalance. The vocabulary (minimum frequency 2, at most 20k types, \texttt{<pad>}$=0$, \texttt{<unk>}$=1$) is built on
the training subset only and contains \Vocab{} tokens.

%-----------------------------------------------------------------------------------------------
\section{Modeling approach}
\paragraph{Architecture.} All models share one pipeline: token ids $\rightarrow$ embedding (128-d) $\rightarrow$ one
recurrent layer built from the chosen cell (hidden size $H=256$) $\rightarrow$ dropout 0.3 $\rightarrow$ linear layer
$\rightarrow$ 2 logits. A generic \texttt{RecurrentLayer} unrolls any cell over time and supports stacking with
inter-layer dropout. Batches are padded to their longest review; at padded positions the state $(\bm h,\bm c)$ is frozen
by a mask, so the classifier receives the hidden state at each review's \textbf{true last token}. For speed, the gate
pre-activations are fused into one \texttt{nn.Linear} for $\bm x_t$ and one bias-free \texttt{nn.Linear} for
$\bm h_{t-1}$ (each gate keeps exactly one bias, as in the equations), and $\bm W_x\bm x_t$ is computed for all time
steps in a single matrix product before the time loop.

\paragraph{Vanilla LSTM (baseline).}
\begin{align*}
\bm i_t &= \sigma(\bm W_{hi}\bm h_{t-1} + \bm W_{xi}\bm x_t + \bm b_i), &
\bm f_t &= \sigma(\bm W_{hf}\bm h_{t-1} + \bm W_{xf}\bm x_t + \bm b_f), &
\bm o_t &= \sigma(\bm W_{ho}\bm h_{t-1} + \bm W_{xo}\bm x_t + \bm b_o),\\
\bm g_t &= \tanh(\bm W_{hg}\bm h_{t-1} + \bm W_{xg}\bm x_t + \bm b_g), &
\bm c_t &= \bm f_t\odot\bm c_{t-1} + \bm i_t\odot\bm g_t, &
\bm h_t &= \bm o_t\odot\tanh(\bm c_t).
\end{align*}

\paragraph{CIFG-LSTM.} The input gate is tied to the forget gate, $\bm i_t = \bm 1 - \bm f_t$; everything else is
unchanged. The cell can only write new content in proportion to what it forgets. Greff et al.\ found that this coupling
costs little accuracy, while it removes a whole gate: \textbf{25\% fewer recurrent parameters} (3 gate blocks instead
of 4).

\paragraph{LSTMP.} The gates are computed from $\bm x_t$ and the \emph{projected} previous state
$\bm h_{t-1}\in\mathbb R^{P}$, while the cell state stays in $\mathbb R^{H}$:
\[
\bm m_t = \bm o_t\odot\tanh(\bm c_t), \qquad \bm h_t = \bm W_p\bm m_t + \bm b_p, \qquad \bm W_p\in\mathbb R^{P\times H},\; P<H .
\]
The four recurrent matrices now act on a $P$-dimensional state, so their cost drops from $4H^2$ to $4HP$ at the price
of one $H\times P$ projection. Sak et al.\ introduced this cell to make large acoustic models cheaper. With $H=256$ and
$P=128$, LSTMP keeps a 256-d memory with roughly the parameter budget of CIFG, which makes the two variants directly
comparable. Both variants therefore ask the same question: \emph{can we keep the vanilla accuracy with fewer
parameters?}

\begin{table}[h]
\centering\small
\ParamsTable
\caption{Trainable parameters ($D=128$, $H=256$). The recurrent counts follow $4H(D+H+1)$ (vanilla), $3H(D+H+1)$
(CIFG) and $4H(D+P+1)+PH+P$ (LSTMP); unit tests check these formulae. Totals include the \Vocab{}$\times$128
embedding, which dominates.}
\label{tab:params}
\end{table}

\paragraph{Verification.} \texttt{tests/test\_cells.py} checks the output shapes, the parameter formulae,
$\bm i_t=\bm 1-\bm f_t$ for CIFG, agreement of the vanilla cell with a direct evaluation of its equations, invariance
of the final state to extra padding, and the absence of built-in recurrent modules in \texttt{src/}.

%-----------------------------------------------------------------------------------------------
\section{Hyperparameter design choices}
\paragraph{Shared protocol.} Every run uses the same splits, the same seed (42), the same seeded length-bucketed batch
order, cross-entropy loss, gradient-norm clipping at 1.0, and at most 10 epochs with early stopping on validation
macro-F1 (patience 3). The checkpoint with the best validation macro-F1 is evaluated \emph{once} on the test set.
Forget-gate biases start at 1 (other biases at 0) so that the cells remember by default early in training.

\paragraph{Baseline.} Adam, learning rate $10^{-3}$, batch size 128, Xavier-uniform initialization, embedding 128,
hidden 256, dropout 0.3. These are standard settings for LSTM text classifiers; the layer sizes were kept moderate so
that each run fits in a free Colab GPU session.

\paragraph{Study.} Starting from the baseline we change \textbf{one factor at a time} on the vanilla model: learning
rate $10^{-3}$ vs $5\cdot10^{-3}$; Adam vs SGD (momentum 0.9, lr 0.05); batch size 64 / 128 / 256; Xavier-uniform vs
orthogonal recurrent weights (Table~\ref{tab:hparams}, Fig.~\ref{fig:hparams}).

\begin{table}[h]
\centering\small
\HparamTable
\caption{Hyperparameter study on the vanilla LSTM (test set; checkpoint chosen on validation macro-F1). All variants have the same 2,783,746 trainable parameters (394,240 recurrent).}
\label{tab:hparams}
\end{table}

\begin{figure}[h]
\centering\figbox{hparam_comparison.png}{0.62\linewidth}
\caption{Test macro-F1 of each hyperparameter variant; the dashed line marks the baseline.}
\label{fig:hparams}
\end{figure}

\HparamFindings

%-----------------------------------------------------------------------------------------------
\section{Final model performance and evaluation}
\paragraph{Architecture comparison.} Table~\ref{tab:arch} compares the three cells with identical hyperparameters.

\begin{table}[h]
\centering\small
\ArchTable
\caption{Architecture comparison on the test set (same splits, seed, hyperparameters and protocol).}
\label{tab:arch}
\end{table}

\ArchFindings

\begin{figure}[h]
\centering
\figbox{val_curves.png}{0.72\linewidth}
\caption{Validation loss and macro-F1 per epoch for the three cells.}
\label{fig:curves}
\end{figure}

\begin{figure}[h]
\centering
\figbox{params_vs_f1.png}{0.72\linewidth}
\caption{Recurrent-layer parameters versus test macro-F1 (the F1 axis does not start at zero).}
\label{fig:params}
\end{figure}

\ProjFindings

\paragraph{Final model.} \BestSummary

\begin{figure}[h]
\centering\figbox{confusion_matrix.png}{0.3\linewidth}
\caption{Test confusion matrix of the final model (counts and row-normalized rates).}
\label{fig:cm}
\end{figure}

\CmSentence

%-----------------------------------------------------------------------------------------------
\section{Limitations and potential improvements}
\begin{itemize}
\item \textbf{Single seed.} Every result comes from one seed, so differences of a few thousandths of macro-F1 are within
noise; repeating each run with 3--5 seeds would give confidence intervals.
\item \textbf{Last-state bottleneck.} A bidirectional layer or attention pooling over all hidden states would avoid
relying on the final state alone, which hurts long reviews whose verdict appears early.
\item \textbf{Pretrained embeddings} (GloVe, fastText, or embeddings trained on medical text) would help with rare
drug names and side-effect vocabulary.
\item \textbf{Other cells.} The Working Memory Connections (WMC) cell, in which the cell state drives the gates, is a
natural fourth variant; stacked layers are already supported by the code.
\item \textbf{Class imbalance.} A class-weighted loss or a decision threshold tuned on validation could raise recall on
the negative class without resampling.
\item \textbf{Task and data.} A 3-class setting that keeps neutral reviews, training on the full 132k training reviews,
and a joint learning-rate/batch-size search would make the conclusions stronger.
\end{itemize}

\vspace{-0.3em}
\paragraph{References.}
{\small
K.~Greff et al., ``LSTM: A Search Space Odyssey,'' \emph{IEEE TNNLS}, 2017 (arXiv:1503.04069).
H.~Sak, A.~Senior, F.~Beaufays, ``Long Short-Term Memory Recurrent Neural Network Architectures for Large Scale
Acoustic Modeling,'' \emph{Interspeech}, 2014.
F.~Landi et al., ``Working Memory Connections for LSTM,'' \emph{Neural Networks}, 2021.
F.~Gr\"a{\ss}er et al., ``Aspect-Based Sentiment Analysis of Drug Reviews Applying Cross-Domain and Cross-Data
Learning,'' \emph{Digital Health}, 2018 (dataset).\par}

\end{document}
